% pdf-style-supplement.tex — extra Unicode->glyph maps for the TPV26/27 plan.
% Loaded AFTER _assets/pdf-style.tex (which already maps · θ μ τ σ Δ ± ∓ ° ⊥ ⟂
% ∈ ≈ ≡ ≤ ≥ ≠ ≪ ≫ × − → ↔ ⇒ ∞ √ ✓ ✗ ⁰ ⁻ ⁿ ⚠).  This file adds the Greek
% letters, super/subscripts, vulgar fractions, integral, section sign and the
% ž that the plan uses but the base header does not cover.
%
% Why \AtBeginDocument: pandoc's xelatex template loads `unicode-math`, which at
% \begin{document} makes several of these codepoints (notably the vulgar
% fractions ½ ⅓ ⅔ ¼, and the super/subscripts) catcode-ACTIVE math characters,
% clobbering a plain preamble \newunicodechar and causing "Missing $ inserted"
% when they appear in text (prose OR inside \texttt code spans).  Declaring our
% maps in \AtBeginDocument — after unicode-math is loaded — runs them AFTER
% unicode-math's own \begin{document} hook, so ours win.  (The base header's
% Greek/relation maps survive because unicode-math sets only their mathcode, not
% catcode, so they don't need this treatment — but wrapping everything is safe.)
\usepackage{newunicodechar}   % harmless no-op if already loaded

\AtBeginDocument{%
  % ---- lowercase Greek ----
  \newunicodechar{α}{\ensuremath{\alpha}}%
  \newunicodechar{β}{\ensuremath{\beta}}%
  \newunicodechar{γ}{\ensuremath{\gamma}}%
  \newunicodechar{δ}{\ensuremath{\delta}}%
  \newunicodechar{ε}{\ensuremath{\varepsilon}}%
  \newunicodechar{η}{\ensuremath{\eta}}%
  \newunicodechar{λ}{\ensuremath{\lambda}}%
  \newunicodechar{ν}{\ensuremath{\nu}}%
  \newunicodechar{ξ}{\ensuremath{\xi}}%
  \newunicodechar{π}{\ensuremath{\pi}}%
  \newunicodechar{ρ}{\ensuremath{\rho}}%
  \newunicodechar{φ}{\ensuremath{\varphi}}%
  \newunicodechar{ω}{\ensuremath{\omega}}%
  \newunicodechar{ψ}{\ensuremath{\psi}}%
  % ---- uppercase Greek ----
  \newunicodechar{Ω}{\ensuremath{\Omega}}%
  \newunicodechar{Σ}{\ensuremath{\Sigma}}%
  \newunicodechar{Φ}{\ensuremath{\Phi}}%
  % ---- superscripts / subscripts (base header maps only ⁰ ⁻ ⁿ) ----
  \newunicodechar{²}{\ensuremath{{}^{2}}}%
  \newunicodechar{³}{\ensuremath{{}^{3}}}%
  \newunicodechar{₀}{\ensuremath{{}_{0}}}%
  \newunicodechar{₁}{\ensuremath{{}_{1}}}%
  \newunicodechar{₂}{\ensuremath{{}_{2}}}%
  \newunicodechar{₃}{\ensuremath{{}_{3}}}%
  \newunicodechar{⁸}{\ensuremath{{}^{8}}}%
  % ---- vulgar fractions: plain ASCII text (unicode-math makes these active;
  %      and \tfrac/\genfrac is fragile inside \texttt / code blocks anyway) ----
  \newunicodechar{½}{1/2}%
  \newunicodechar{⅓}{1/3}%
  \newunicodechar{⅔}{2/3}%
  \newunicodechar{¼}{1/4}%
  % NB: use \intop, NOT \int — under unicode-math \int is \let to the active
  % char ∫, so \newunicodechar{∫}{...\int...} recurses infinitely (a real hang).
  \newunicodechar{∫}{\ensuremath{\intop}}%
  % ---- misc text glyphs ----
  \newunicodechar{§}{\S}%
  \newunicodechar{←}{\ensuremath{\leftarrow}}%
  \newunicodechar{ž}{\v{z}}%
}

% ---- Robust code-box: override pdf-style.tex's breakable tcolorbox `Shaded`.
% This plan has long fenced code blocks that must break across pages; the
% breakable tcolorbox loops on "Infinite glue shrinkage found in box being
% split" and never finishes.  framed's `snugshade` (already loaded via
% \usepackage{framed}, and pandoc's own default) page-breaks cleanly.  A second
% \AtBeginDocument declared AFTER pdf-style.tex's runs later, so this wins. ----
\usepackage{framed}
\definecolor{shadecolor}{HTML}{F6F8FA}
\AtBeginDocument{%
  \ifcsname Shaded\endcsname
    \renewenvironment{Shaded}{\begin{snugshade}}{\end{snugshade}}%
  \fi
}

% ---- Wrap long lines inside code blocks instead of overflowing the margin.
% Plan docs are full of 100+ char repo paths and CLI invocations; without this
% every such line is an Overfull \hbox that runs off the page.  fvextra's
% breaklines/breakanywhere covers BOTH pandoc's highlighted `Highlighting`
% environment (a fancyvrb Verbatim) and untagged ``` fences (plain `verbatim`).
% Must be loaded AFTER pandoc's own \DefineVerbatimEnvironment{Highlighting},
% which the template emits near the top of the preamble — any -H header
% qualifies.  The continuation hook (fvextra's default breaksymbolleft) marks
% wrapped lines so a broken path is not mistaken for two paths.
\usepackage{fvextra}
\DefineVerbatimEnvironment{Highlighting}{Verbatim}%
  {breaklines,breakanywhere,commandchars=\\\{\}}
\RecustomVerbatimEnvironment{verbatim}{Verbatim}%
  {breaklines,breakanywhere}
